\chapter{Bacterial Infections}

\section*{Definition and Overview}
Bacterial infections are illnesses caused by bacteria -- single-celled micro-organisms that can live and multiply on their own. Unlike viruses, bacteria are complete cells that can survive and reproduce in many environments, including on and inside the human body. Many bacteria are harmless or even helpful, but some cause disease by damaging tissues or releasing toxins.

Bacterial infections can affect almost any part of the body, including the skin, throat, chest, urinary tract, gut, bones, and bloodstream. They range from minor, self-limiting infections such as a mild skin boil to life-threatening conditions such as pneumonia, meningitis, or sepsis. The key feature that separates bacterial from viral infection is that bacteria are killed by antibiotics, which is why antibiotics are prescribed for bacterial infections but are ineffective against viruses.

\section*{Epidemiology}
Bacterial infections are among the most common reasons people see a doctor. Everyday examples include throat infections (strep throat), ear infections, urinary tract infections, and skin infections such as impetigo and cellulitis. Some serious bacterial diseases -- such as tuberculosis, meningococcal disease, and whooping cough -- occur worldwide and are more common where vaccination, sanitation, and healthcare access are limited.

Antibiotic resistance is a growing problem: bacteria such as methicillin-resistant \textit{Staphylococcus aureus} (MRSA) and resistant tuberculosis are harder to treat. Vaccination has greatly reduced diseases caused by bacteria such as \textit{Haemophilus influenzae} type b, \textit{Streptococcus pneumoniae}, and \textit{Neisseria meningitidis}. People at the extremes of age, the immunocompromised, and those in crowded or resource-poor settings carry the highest burden of severe disease.

\section*{Aetiology and Pathophysiology}
Bacteria cause disease in a few ways, and the site of infection determines the symptoms:

\begin{itemize}
    \item \textbf{Tissue damage:} bacteria multiply locally and directly injure cells (for example, \textit{Streptococcus pyogenes} in the throat, \textit{Escherichia coli} in the urinary tract)
    \item \textbf{Toxin production:} some bacteria release poisons, such as the toxins that cause tetanus, botulism, or the diarrhoea of certain food poisoning
    \item \textbf{Immune response:} the body's own inflammatory response to the bacteria can cause much of the damage, including fever, swelling, and, in severe cases, sepsis
    \item \textbf{Transmission:} spread by droplets (coughs and sneezes), touch and skin contact, contaminated food and water, sexual contact, or through breaks in the skin
    \item \textbf{Common examples:} \textit{Staphylococcus aureus} (skin), \textit{Streptococcus} (throat), \textit{E. coli} (urine and gut), \textit{Salmonella} (food poisoning), \textit{Mycobacterium tuberculosis} (tuberculosis), and \textit{Neisseria meningitidis} (meningitis)
\end{itemize}

\section*{Clinical Features}
Symptoms depend on the part of the body infected, but common features include:

\begin{itemize}
    \item Fever, chills, and feeling generally unwell
    \item Localised pain, redness, swelling, and warmth (for example, skin infection, wound infection)
    \item Pus or discharge from a wound, throat, ear, or other site
    \item \textbf{Respiratory infection:} sore throat, cough with coloured phlegm, chest pain, shortness of breath
    \item \textbf{Urinary tract infection:} burning with urination, urgency, cloudy or smelly urine
    \item \textbf{Gastrointestinal infection:} diarrhoea, vomiting, abdominal cramps
    \item \textbf{Meningitis:} severe headache, stiff neck, dislike of bright light, sometimes a rash
    \item \textbf{Sepsis:} rapid breathing, rapid heart rate, confusion, mottled or cold skin, reduced urine output
\end{itemize}

\section*{Differential Diagnosis}
\begin{itemize}
    \item Viral infection -- the most common mimic (for example, viral versus bacterial sore throat or pneumonia)
    \item Fungal infection, which can look similar on the skin or in the lungs
    \item Non-infectious inflammation (for example, autoimmune conditions, gout, allergy)
    \item Malignancy, which can present with fever, sweats, and weight loss
    \item Parasitic infections, in relevant settings (for example, malaria)
    \item Non-infectious fever: drug reactions, blood clots, or heat illness
\end{itemize}

\section*{When to Seek Medical Care}
Mild bacterial infections often need treatment, so see a doctor if you suspect one -- for example, a sore throat with high fever and swollen glands, a spreading skin infection, burning on urination, or a cough with coloured phlegm that is not improving.

\textbf{Call 000 (or local emergency services) for:}
\begin{itemize}
    \item Difficulty breathing, rapid breathing, or chest pain
    \item Confusion, drowsiness, or a stiff neck with a rash
    \item A spreading rash that does not blanch (fade when pressed)
    \item Signs of sepsis: mottled or cold skin, very fast heart rate, low urine output, dizziness or fainting
    \item High fever that does not respond to treatment, especially in babies, older adults, or people with weakened immunity
\end{itemize}

\section*{Diagnosis and Investigations}
\begin{itemize}
    \item \textbf{History and examination:} site of infection, onset, exposures, fever, and risk factors such as diabetes or immune suppression
    \item \textbf{Blood tests:} full blood count (raised white cells), inflammatory markers such as C-reactive protein, and blood cultures in serious illness
    \item \textbf{Specimen testing:} throat swab, urine culture, wound swab, or stool culture to identify the bacteria and test antibiotic sensitivity
    \item \textbf{Molecular tests:} PCR can quickly detect specific bacteria, such as \textit{Chlamydia}, pertussis, or meningococcus
    \item \textbf{Imaging:} chest X-ray for pneumonia; CT or ultrasound to find a hidden focus such as an abscess
    \item \textbf{Lumbar puncture:} testing of spinal fluid when meningitis is suspected
\end{itemize}

\section*{Management}
\begin{itemize}
    \item \textbf{Antibiotics:} the main treatment; chosen against the likely or confirmed bacteria, and adjusted once sensitivity results are known
    \item \textbf{Take the full course:} even when symptoms improve, finishing the prescribed course helps prevent resistance and relapse
    \item \textbf{Supportive care:} rest, fluids, and simple pain and fever relief such as paracetamol
    \item \textbf{Drainage:} abscesses, boils, and infected fluid collections usually need draining as well as antibiotics
    \item \textbf{Source control:} removal of infected devices, cleaning of wounds, or surgery for deep or necrotising infections
    \item \textbf{Hospital care:} intravenous antibiotics, fluids, and oxygen for severe infections and sepsis
    \item \textbf{Prevention of spread:} good hand hygiene and staying home while infectious
\end{itemize}

\section*{Complications}
\begin{itemize}
    \item Sepsis and septic shock -- a life-threatening whole-body response to infection
    \item Abscess formation, which may need drainage
    \item Spread of infection to the bloodstream, bones, joints, or heart valves
    \item Organ damage, including kidney, lung, or heart complications
    \item Post-infectious conditions, such as rheumatic fever after untreated strep throat
    \item Antibiotic resistance and recurrent infections
    \item Chronic infection, for example tuberculosis, if not fully treated
\end{itemize}

\section*{Prognosis and Outcomes}
Most bacterial infections resolve completely with appropriate antibiotics and supportive care. Simple infections such as strep throat and urinary tract infections usually improve within a couple of days of starting the right antibiotic. Serious infections such as pneumonia, meningitis, and sepsis carry significant risk, and outcomes depend on the bacteria involved, how quickly treatment starts, and the person's age and general health. Untreated bacterial infections can progress, which is why prompt assessment matters. With correct treatment, most people make a full recovery.

\section*{Prevention and Self-Care}
\begin{itemize}
    \item Keep vaccinations up to date, including whooping cough, pneumococcal, and meningococcal vaccines
    \item Wash hands regularly with soap and water, especially after using the toilet and before eating
    \item Prepare and store food safely; cook meat thoroughly and refrigerate leftovers promptly
    \item Clean and cover cuts and wounds, and keep them dry
    \item Do not share towels, razors, or drink bottles
    \item Use condoms to reduce sexually transmitted bacterial infections
    \item Take antibiotics exactly as prescribed and never share them
    \item Stay home when unwell to avoid spreading infection to others
\end{itemize}

\section*{Sources and Further Reading}
The following sources provide reliable, up-to-date information on bacterial infections.
\begin{itemize}
    \item Centers for Disease Control and Prevention. \textit{Bacterial infections.} https://www.cdc.gov
    \item World Health Organization. \textit{Antimicrobial resistance.} https://www.who.int
    \item NHS. \textit{Infections.} https://www.nhs.uk/conditions/
    \item MSD Manual. \textit{Overview of Bacteria.} https://www.msdmanuals.com
    \item MedlinePlus. \textit{Bacterial Infections.} https://medlineplus.gov
\end{itemize}
